
\section{The number of firing rates grows at least like $A^{2/3}$}\label{sec:thmC}

Let $K(A,\lambda)$ be the number of mass points of the capacity-achieving input of the channel $Y\sim\mathrm{Poisson}(X)$ with
$\lambda\le X\le\lambda+A$, where $\lambda\ge0$ is the spontaneous count, and $K(A)=K(A,0)$. The constant of the theorem comes from a one-dimensional optimization over a parameter
$\xi$ of the proof; a reader interested only in the result may skip the definitions that follow. In words, $c_P$ below is the best
constant that the construction of the proof yields, and $\lambda_M$ and $\lambda_\ell$ are numbers of that construction, unrelated to
the spontaneous count $\lambda$. For $\xi\in(0,1)$ put $\phi_\xi=\xi\log\xi+1-\xi$, let $\kappa_L(\xi)\in(0,\xi)$ be the root of
$\log k+(1-k)/\xi=0$, and put
\[
h(f)=\frac{\xi\log f+1-f}{\phi_\xi},\qquad
\mathcal L(1,z)=\frac{z-1}{\log z},\qquad
\varrho=\frac{\mathcal L(1,M)}{\mathcal L(1,\ell)},
\]
where $M=e^{\lambda_M}$ with $\lambda_M=1-\psi(2)+e^{\psi(2)}$, $\ell=e^{-\lambda_\ell}$ with $\lambda_\ell=2.777$, and $\psi$ is the
digamma function. Let
\[
c_{0}^{\max}(\xi)=\Bigl[\bigl(1+\xi\kappa_L^{-1/2}\bigr)\kappa_L^{-1/4}\varrho^{1/2}/\phi_\xi\Bigr]^{-2/3},\qquad
\mathcal I(\xi)=\int_{\kappa_L}^{1}\Bigl(\frac{h(f)}{f}\Bigr)^{1/2}df .
\]
Both quantities enter the bound as follows.

\begin{theorem}\label{thm:C}
For every fixed $\lambda\ge0$,
\[
\liminf_{A\to\infty}\frac{K(A,\lambda)}{A^{2/3}}\;\ge\;c_P:=\max_{0<\xi<1}\frac{c_{0}^{\max}(\xi)^{1/2}\,\mathcal I(\xi)}{2\pi}\;>\;0.0560 .
\]
At $\xi=0.2411$, near the maximizer, $c_{0}^{\max}=0.06491$, $\mathcal I=1.38249$ and the bound is $0.05606$.
\end{theorem}

In words: for every $\varepsilon>0$ and every large enough $A$, the optimal input has at least $(c_P-\varepsilon)\,A^{2/3}$ mass
points; the bound concerns growth, and the lower limit $\liminf$ allows finitely many exceptions at small $A$. The proof is given
in Appendix~\ref{app:C}; its constants are enclosed in interval arithmetic, so the numerical steps of the proof are proved bounds. In the Fisher half-length $a=\sqrt A$ the bound reads $K\ge(c_P-o(1))\,a^{4/3}$, where $o(1)$ tends to zero as $A$ grows. This is the exponent of the count
law of the Gaussian channel: the companion paper~\cite{Paper1} proves that the Gaussian channel with input in $[-a,a]$ has at least
$0.166\,a^{4/3}$ mass points for large $a$, and measures about $0.4002\,a^{4/3}$. The best published lower bound on the same number $K(A)$
that we know of is of order $\sqrt A$~\cite{DBS2021,BDS2024}, so the theorem raises the proved exponent from $1/2$ to $2/3$. The
best published upper bound is of order $A$~\cite{BDS2024}, so the proved order of growth now lies between $A^{2/3}$ and $A$; the law
of Section~\ref{sec:countlaw} places it at $A^{2/3}$. The constant of Theorem~\ref{thm:C} is far from sharp. At the peak count
$2000$ the bound gives about 9 steps, while the optimal input has 66; the theorem gives the order of growth from below and leaves the
exact count to the law of Section~\ref{sec:countlaw}.

\subsection{Setting and facts from the literature}

Until Section~\ref{app:C-lam} the spontaneous count is $\lambda=0$; that section gives the changes for $\lambda>0$.
For $y\in\mathbb N_0=\{0,1,2,\dots\}$ put $p(y\,|\,x)=e^{-x}x^y/y!$; $P_x$ and $E_x$ denote probability and expectation for $Y\sim\mathrm{Poisson}(x)$. Let $\pi=\sum_{j=1}^{K}w_j\delta_{x_j}$, with weights $w_j$, be the capacity-achieving input, $m=\mathsf P\pi$ its
output law ($\mathsf P$ is the Poisson kernel), $C$ the capacity in nats and $D(x)=D(p(\cdot\,|\,x)\,\|\,m)$.
\begin{itemize}
\item[(F1)] $\pi$ exists, is unique and has finitely many mass points~\cite{Shamai1990}; it is optimal if and only if $D(x)\le C$ on
$[0,A]$, with equality on its support~\cite[Lemma~1]{BDS2024}.
\item[(F2)] $0$ and $A$ are mass points of $\pi$~\cite[Corollaries~2 and~4]{Cao2014}.
\item[(F3)] $E[X\,|\,Y=y]=(y+1)m(y+1)/m(y)$ and $\frac{d}{dx}E[f(Y)\,|\,X=x]=E[f(Y+1)-f(Y)\,|\,X=x]$ for $x>0$~\cite[Appendix~D]{BDS2024}.
Since $p(y_0\,|\,x)=p(y_0-1\,|\,x)\,x/y_0$, the posterior of $X$ given $Y=y_0$ is the posterior given $Y=y_0-1$ tilted by $x$.
\item[(F4)] $E[X\,|\,Y=y]$ is nondecreasing in $y$~\cite[Lemma~3]{BDS2024}. With (F3) this is the convexity of
$G(y)=\log(y!\,m(y))$, whose increments are $\log E[X\,|\,Y=y]$.
\end{itemize}

\subsection{A reference output and the exact form of the optimality conditions}

Let $\varphi(y)=y\psi(y+1)-y$, so that $\varphi(y+1)-\varphi(y)=\psi(y+1)$, and put
\[
r(y)=\frac{e^{-(C+1)+\varphi(y)}}{y!},\qquad c_*=r(0)=e^{-(C+1)},\qquad g=\log(m/r).
\]
The function $r$ is a reference output law, compared with $m$ below.

\begin{lemma}\label{lem:C-digamma}
For $x\ge0$ one has $E_x\psi(Y+1)=\log x+E_1(x)$ and $E_x\varphi(Y)=x\log x-x+1-\delta_0(x)$, where $\delta_0(x)=e^{-x}-xE_1(x)$ satisfies
$0<\delta_0(x)\le e^{-x}/(x+1)$ and $\delta_0(0)=1$.
\end{lemma}
\begin{proof}
With Euler's constant $\gamma_E$ one has $\psi(y+1)=-\gamma_E+H_y$ with $H_y=\int_0^1(1-(1-s)^y)\,s^{-1}ds$, so $E_xH_Y=\int_0^1(1-e^{-xs})s^{-1}ds=\gamma_E+\log x+E_1(x)$.
By (F3), $\frac{d}{dx}E_x\varphi(Y)=E_x\psi(Y+1)$, and $E_0\varphi(Y)=0$; integrating gives
$\int_0^x(\log s+E_1(s))\,ds=x\log x-x+xE_1(x)-e^{-x}+1$. The bounds follow from $e^{-x}/(x+1)<E_1(x)<e^{-x}/x$.
\end{proof}

\begin{proposition}\label{prop:C-kkt}
For every $x\ge0$ one has $D(p(\cdot\,|\,x)\,\|\,r)=C+\delta_0(x)$. Hence $C-D(x)=E_xg(Y)-\delta_0(x)$, and the optimality conditions read
$E_xg(Y)\ge\delta_0(x)$ on $[0,A]$, with equality at every mass point. We also have $g(0)=1$.
\end{proposition}
\begin{proof}
The log ratio is $\log p(y\,|\,x)-\log r(y)=-x+y\log x+C+1-\varphi(y)$; take $E_x$ and use Lemma~\ref{lem:C-digamma} with $E_xY=x$. Since $0$ is a
mass point (F2) and $p(\cdot\,|\,0)=\delta_0$, $C=D(0)=-\log m(0)$, so $m(0)=e^{-C}=e\,r(0)$.
\end{proof}

For $z=m/r>0$ one has $(z-1)\log z=\mathcal L(1,z)(\log z)^2$ and $(z-1)^2=\mathcal L(1,z)(z-1)\log z$; so $g$ has the sign of $m-r$ and
$(m-r)g\ge0$ pointwise.

\begin{proposition}[weighted identity]\label{prop:C-weighted}
For every bounded measurable $\chi\ge0$ on $[0,A]$ and every finite measure $\rho\ge0$ on $[0,A]$, with $\mu=\rho-\pi$,
\[
\mathcal W(\chi):=\int\chi\,(C-D)\,d\rho\ \ge0,\qquad
\mathcal W(\chi)=\sum_{y\ge0}g(y)\,\mathsf P(\chi\mu)(y)-\int\chi\,\delta_0\,d\mu .
\]
\end{proposition}
\begin{proof}
One has $\mathcal W\ge0$ by the optimality inequality, and $\int\chi(C-D)\,d\pi=0$ by the equality on the support. Insert
Proposition~\ref{prop:C-kkt} and exchange sum and integral; the exchange is justified because $|g(y)|$ grows at most quadratically
in $y-\lfloor A\rfloor$ beyond $A$ (Lemma~\ref{lem:C-edge}(iii)), while $\mathsf P(\chi|\mu|)(y)$ decays like $p(y\,|\,A)$ for $y>A$.
\end{proof}

\subsection{The reference input}\label{app:C-ref-sec}

\begin{lemma}\label{lem:C-ref}
(i) Let $\pi_Z(u)=(ue^u-e^u+1)/(u(e^u-1)^2)>0$ on $(0,\infty)$, let $S\ge0$ be the driftless subordinator variable (the value at time one of an increasing L\'evy process without drift; only this random variable is used, never the process) with Laplace
exponent $\int(1-e^{-pu})\pi_Z(u)\,du$, that is, $E\,e^{-pS}=\exp\bigl(-\int(1-e^{-pu})\pi_Z(u)\,du\bigr)$, and put $Z=e^{-S}$. Then $EZ^p=\exp(p\psi(p+1)-p-\log\Gamma(p+1))$ for $p\ge0$, so
$r(y)=c_*\,EZ^y$.
(ii) $\rho_\infty(dx)=c_*\,E[Z^{-1}e^{-x(1-Z)/Z}]\,dx$ is a completely monotone density on $(0,\infty)$ with $\mathsf P\rho_\infty=r$.
(iii) $\hat\rho:=\rho_\infty/c_*$ satisfies $\hat\rho(x)\le x^{-1/2}$ for $x\ge1$.
(iv) The restriction $\rho_A$ of $\rho_\infty$ to $[0,A]$ has mass $m_A\le\exp(-\int\delta_0\,d\rho_A/m_A)<1$.
\end{lemma}
\begin{proof}
(i) $k(p):=p\psi(p+1)-p-\log\Gamma(p+1)$ has $k(0)=0$ and $k'(p)=p\psi'(p+1)-1$. From
$\psi'(p+1)=\int_0^\infty ue^{-pu}(e^u-1)^{-1}du$, integration by parts gives $p\psi'(p+1)=1+\int e^{-pu}\frac{d}{du}[u/(e^u-1)]\,du$, and
$\frac{d}{du}[u/(e^u-1)]=-u\pi_Z(u)$. Hence $k(p)=-\int(1-e^{-pu})\pi_Z(u)\,du$; since $\int\min(1,u)\pi_Z<\infty$, this is the Laplace
exponent of a driftless subordinator. We also have $\pi_Z>0$ because $ue^u-e^u+1>0$ for $u>0$.
(ii) The density $z^{-1}e^{-x(1-z)/z}$ on $(0,\infty)$ has Poisson output $z^y$, because $\int_0^\infty e^{-x/z}x^y/y!\,dx=z^{y+1}$;
mix over $z=Z$ (positive terms, Fubini).
(iii) $\pi_Z(u)\ge e^{-u}/(2u)$: with $v=e^u$ this is $2v^2\log v-3v^2+4v-1\ge0$ for $v\ge1$, which vanishes at $v=1$ with zero
derivative and has second derivative $4\log v\ge0$. The right side is the L\'evy density of the $\mathrm{Gamma}(1/2,1)$ subordinator,
so $S=S_1+S'$ with $S_1$ a $\mathrm{Gamma}(1/2,1)$ variable and $S'\ge0$ independent of $S_1$. For $x\ge1$, $s\mapsto e^{s-x(e^s-1)}$ is decreasing, and
$\hat\rho(x)=E\,e^{S-x(e^S-1)}\le E\,e^{S_1-x(e^{S_1}-1)}\le E\,e^{-(x-1)S_1}=x^{-1/2}$.
(iv) $\bar\rho=\rho_A/m_A$ is a feasible input with output $R_A/m_A$, where $R_A=\mathsf P\rho_A\le r$. By Proposition~\ref{prop:C-kkt},
$I(\bar\rho)=C+\int\delta_0\,d\bar\rho+\log m_A-\sum(R_A/m_A)\log(R_A/r)\ge C+\int\delta_0\,d\bar\rho+\log m_A$, and $I(\bar\rho)\le C$.
\end{proof}

\begin{lemma}[two-sided sizes]\label{lem:C-sizes}
(i) $\hat\rho(x)\ge c_1/(x\log x)^{1/2}$ for $x\ge x_0$, with absolute constants $c_1,x_0$.
(ii) Under Proposition~\ref{prop:C-lower}, $c_*\le C_1A^{-1/2}$; with Lemma~\ref{lem:C-edge}(vii), $c_*=\Theta(A^{-1/2})$.
\end{lemma}
\begin{proof}
(i) $e^{s-x(e^s-1)}\ge e^{-2xs}$ for $0\le s\le1$, so $\hat\rho(x)\ge e^{-2}P(S\le1/x)$. Since $Ee^{-pS}\le P(S\le s)+e^{-ps}$ and
$Ee^{-pS}=r(p)/c_*\ge(e/(2\pi p))^{1/2}e^{-5/(24p)}$ for $p\ge1$ (bounds on $r$ in Section~\ref{sec:C-range}), the choice $s=1/x$,
$p=2x\log x$ gives $P(S\le1/x)\ge c/(x\log x)^{1/2}-x^{-2}\ge c'/(x\log x)^{1/2}$ for large $x$.
(ii) $1\ge\sum_{y\le A}m(y)\ge\ell\sum_{y\le A}r(y)\ge\ell c_*\sum_{1\le y\le A}(e/(2\pi y))^{1/2}e^{-5/(24y)}\ge c''\ell c_*A^{1/2}$.
\end{proof}

The right-edge remainder $T_A=r-R_A=\mathsf P(\rho_\infty\mathbf 1_{(A,\infty)})$ satisfies $T_A(y)\le c_*A^{-1/2}e^{-(A-y)^2/(2A)}$
for $y\le A$.

\subsection{The range of $m/r$}\label{sec:C-range}

By (F4) $G(y)=\log(y!\,m(y))$ is convex, and $g=G-\varphi+C+1$. For real $y\ge1$,
\[
\sqrt{e/(2\pi y)}\,e^{-5/(24y)}\le r(y)/c_*\le\sqrt{e/(2\pi y)}\,e^{1/(2(y+1))},
\]
and $r$ is decreasing ($r=c_*EZ^y$ with $Z<1$); this follows from $\log(y+\tfrac12)<\psi(y+1)<\log(y+1)-1/(2(y+1))$ and Stirling's
bounds.

\begin{lemma}[pointwise upper bound]\label{lem:C-pointwise}
For $y_0\ge1$ let $m_1=E[X\,|\,Y=y_0-1]$ and $m_2=E[X\,|\,Y=y_0]$. Then
$g(y_0)\le1+y_0\log m_1-m_1-\varphi(y_0)+m_1\log(m_2/m_1)$.
\end{lemma}
\begin{proof}
The slope $s_0=G(y_0)-G(y_0-1)=\log m_1$ (F3) is a subgradient of the convex sequence $G$ at $y_0$, so
$g(y)\ge g(y_0)+s_0(y-y_0)-(\varphi(y)-\varphi(y_0))$ for all $y\in\mathbb N_0$. Take $E_{x_j}$ at a mass point $x_j$:
$E_{x_j}g(Y)=\delta_0(x_j)$ (Proposition~\ref{prop:C-kkt}) and $E_{x_j}\varphi(Y)=x_j\log x_j-x_j+1-\delta_0(x_j)$
(Lemma~\ref{lem:C-digamma}), so $\delta_0$ cancels and $g(y_0)\le1+[s_0y_0-e^{s_0}-\varphi(y_0)]+B(x_j,e^{s_0})$ with
$B(x,x')=x\log(x/x')-x+x'\ge0$. Average over the posterior of $X$ given $Y=y_0-1$, a probability on the mass points with mean
$m_1=e^{s_0}$. The posterior given $Y=y_0$ is that posterior tilted by $x$, so $E[B(X,m_1)\,|\,Y=y_0-1]=m_1E[\log(X/m_1)\,|\,Y=y_0]
\le m_1\log(m_2/m_1)$ by Jensen's inequality.
\end{proof}

\begin{proposition}\label{prop:C-upper}
The function $g$ satisfies $\sup_{y\in\mathbb N_0}g(y)\le\lambda_M=1-\psi(2)+e^{\psi(2)}=2.1034207765\ldots$, that is, $m\le Mr$.
\end{proposition}
\begin{proof}
One has $g\to-\infty$ because $m$ has Poisson tails beyond $A$ while $r(y)$ is of order $y^{-1/2}$, so $g$ attains its maximum at some $y_0$.
If $y_0=0$, $g(0)=1$ (Proposition~\ref{prop:C-kkt}). If $y_0\ge1$, then $g(y_0-1)\le g(y_0)\ge g(y_0+1)$ and $\varphi(y)-\varphi(y-1)=\psi(y)$ give
$m_1\ge e^{\psi(y_0)}$ and $m_2\le e^{\psi(y_0+1)}$. In Lemma~\ref{lem:C-pointwise}, $y_0\log m_1-m_1\le y_0\log y_0-y_0$, and
$m_1\log(m_2/m_1)\le m_1(\psi(y_0+1)-\psi(y_0))=m_1/y_0\le m_2/y_0\le e^{\psi(y_0+1)}/y_0$. With $\varphi(y_0)=y_0\psi(y_0+1)-y_0$,
\[
g(y_0)\le f(y_0):=1+y_0(\log y_0-\psi(y_0+1))+e^{\psi(y_0+1)}/y_0 .
\]
At $y_0=1$ the bound is $f(1)=\lambda_M$. For $y_0\ge2$ the inequalities $\psi(y_0+1)>\log(y_0+\frac12)$ and $\log(1+u)\ge u-u^2/2$ give
$y_0(\log y_0-\psi(y_0+1))\le-\frac12+1/(8y_0)$, and $\psi(y_0+1)<\log(y_0+1)$ gives $e^{\psi(y_0+1)}/y_0<1+1/y_0$; so
$f(y_0)<\frac32+9/(8y_0)\le2.0625<\lambda_M$.
\end{proof}

\begin{proposition}[lower bound on the input range]\label{prop:C-lower}
For every integer $y_1\in[0,A]$ one has $g(y_1)\ge-\lambda_\ell$ with $\lambda_\ell=2.777$, that is, $m\ge\ell r$ there.
\end{proposition}
\begin{proof}
Fix $y_1$ and put $\lambda_1=-g(y_1)$; there is nothing to show if $\lambda_1\le0$. For an integer $R\ge1$ and $y=y_1+j$ with $0\le j\le R$, convexity of $G$ and $g(y_1+R)\le\lambda_M$ give $g(y)\le(1-\theta)(-\lambda_1)+\theta\lambda_M+\mathrm{def}_R(j)$ with $\theta=j/R$ and the
chord defect $\mathrm{def}_R(j)=(1-\theta)\varphi(y_1)+\theta\varphi(y_1+R)-\varphi(y_1+j)\le j(R-j)/(2(y_1+1))$, by the discrete maximum
principle since $\Delta^2\varphi(y)=1/(y+1)$. Left chords ($R\le y_1$) give the same bound with $y_1+1$ replaced by $y_1-R+1$, and beyond
the chords $g\le\lambda_M$. Let $U_{\lambda_1}$ be the pointwise minimum of these bounds; it is nonincreasing in $\lambda_1$. The optimality
inequality at the feasible point $x=y_1$ gives $E_{y_1}U_{\lambda_1}\ge E_{y_1}g\ge\delta_0(y_1)\ge0$.

(a) $y_1\le3000$. For $y_1=0$ the claim holds because $g(0)=1$ (Proposition~\ref{prop:C-kkt}). The bound $U_{1.8}$ is built from chord bounds, with the Poisson mass outside the chords counted at
$\lambda_M$. For each integer $y_1\in[1,3000]$ it gives $E_{y_1}U_{1.8}-\delta_0(y_1)<0$, with an upper bound of at most $-6.67\times10^{-3}$
(attained at $y_1=4$); each value is enclosed in ball arithmetic (interval arithmetic with midpoint-radius balls) at 128 bits. Hence $\lambda_1<1.8$ there.

(b) $y_1>3000$. One chord on each side with $R=\lceil c\sqrt{y_1}\,\rceil$, $c=2.6$, and $\theta_Y=|Y-y_1|/R$ give
$U\le-\lambda_1(1-\theta_Y)_++\lambda_M\min(1,\theta_Y)+R^2/(8(y_1-R+1))$. Since $E|Y-y_1|\le\sqrt{y_1}$, it follows that $E(1-\theta_Y)_+\ge1-1/c$ and
$E\min(1,\theta_Y)\le1/c$; so $\lambda_1\le(\lambda_M/c+R^2/(8(y_1-R+1)))/(1-1/c)$. With $R\le c\sqrt{y_1}+1$ the last term is at most
$R_b(y_1)=(c^2y_1+2c\sqrt{y_1}+1)/(8(y_1-c\sqrt{y_1}))$, which decreases on $y_1\ge3001$, and $R_b(3001)$ gives $\lambda_1\le2.77649$.
\end{proof}

\begin{corollary}\label{cor:C-LM}
With $z=m/r$: (LM1) $(z-1)\log z\ge\mathcal L(1,\ell)(\log z)^2$ on $[0,A]$, and (LM2) $(z-1)^2\le\mathcal L(1,M)(z-1)\log z$ on
$\mathbb N_0$. The value of $\varrho$ is $10.128$.
\end{corollary}
\begin{proof}
The ratio satisfies $z\in[\ell,M]$ on $[0,A]$ and $z\le M$ on $\mathbb N_0$ (Propositions~\ref{prop:C-upper} and~\ref{prop:C-lower}), and
$\mathcal L(1,z)$ is increasing in $z$.
\end{proof}

\subsection{Exponential decay of $m-r$ in the bulk}

\emph{Weight.} Fix $\xi\in(0,1)$ and an integer $\sigma\ge1$, and put $t=2\sigma/(\xi A)$ and $\chi(x)=(x/A)^{2\sigma}e^{t(A-x)}$. Then
$\log\chi(x)=\Lambda_0h(x/A)$ with $\Lambda_0=tA\phi_\xi$; $\log\chi$ is concave, maximal at $x_c=\xi A$, zero at $x_L=\kappa_LA$ and at
$A$, positive on $(x_L,A)$ and negative on $[0,x_L)$ and beyond $A$, with $\chi(0)=0$. Tangent bounds:
$\log\chi(x)\le\beta_L(x-x_L)$ with $\beta_L=t(\xi/\kappa_L-1)$, and $\log\chi(x)\le\beta_A(A-x)$ with $\beta_A=t(1-\xi)$.
Let $\mathrm{Win}=\{y\in\mathbb N_0:x_L\le y\le A\}$.

\emph{Objects.} Put $\rho=\rho_A$ (Lemma~\ref{lem:C-ref}), $\mu=\rho-\pi$, $\nu=\chi^{1/2}\mu$ (which has no mass at the origin since $\chi(0)=0$),
with $N:=\|\nu\|_{\mathrm{TV}}\le2e^{\Lambda_0/2}$. Put $\lambda_\pm=1\pm t/2$ and $H_\lambda(k)=\int p(k\,|\,\lambda x)\,\nu(dx)$;
$F=\mathsf P\mu=r-T_A-m$, $\mathcal P(y)=\chi(y)^{1/2}F(y)$, $\mathcal Q(y)=\chi(y)^{-1/2}\mathsf P(\chi\mu)(y)$ and
$C_m=\mathsf P(\chi\mu)-\chi F=\chi^{1/2}(\mathcal Q-\mathcal P)$. The weighted error is
$E_2=\sum_{y\le A}\chi(y)(m-r)g\ge0$.

\begin{proposition}[decay]\label{thm:C-decay}
Fix $\xi\in(0,1)$ and $c_0<c_{0}^{\max}(\xi)=\Gamma_\xi^{-2/3}$, where
\[\Gamma_\xi=(1+\xi\kappa_L^{-1/2})\kappa_L^{-1/4}\varrho^{1/2}/\phi_\xi.\]
Take $\sigma=\mathrm{round}(c_0A^{1/3}\xi/(2\phi_\xi))$, so that $\Lambda_0=c_0A^{1/3}(1+o(1))$, and $B^2=\lambda_+(\Lambda_0+4\log A)$.
There are $A_0$ and $C_0$, depending on $\xi$ and $c_0$ only, such that $E_2\le C_0A^{1/2}$ for $A\ge A_0$.
\end{proposition}
\begin{proof}
\emph{Step 1 (identity).} Proposition~\ref{prop:C-weighted} with this $\chi$ and $\rho$, and $\mathsf P(\chi\mu)=\chi F+C_m$ on $y\le A$,
give $E_2=-\mathcal W-\tilde T+\sum_{\mathrm{Win}}gC_m+e_L+e_{\rm out}-d_0$, with $\tilde T=\sum_{y\le A}\chi gT_A$ and $e_L=\sum_{y<x_L}gC_m$, and with $e_{\rm out}=\sum_{y>A}g\,\mathsf P(\chi\mu)$ and $d_0=\int\chi\delta_0\,d\mu$. Since $\mathcal W\ge0$, the bound $E_2\le|\tilde T|+|\sum_{\mathrm{Win}}gC_m|+|e_L|+|e_{\rm out}|+|d_0|$ holds.

\emph{Step 2 (exact Poisson flows).} From
$x^se^{-\tau x}p(y\,|\,x)=\frac{\Gamma(y+s+1)}{\Gamma(y+1)(1+\tau)^{y+s}}\,p(y+s\,|\,(1+\tau)x)$ with $s=\pm\sigma$, $\tau=\pm t/2$:
for $y\ge\sigma$ one has $\mathcal Q(y)=\alpha_+(y)H_{\lambda_+}(y+\sigma)$ and $\mathcal P(y)=\alpha_-(y)H_{\lambda_-}(y-\sigma)$, with
$\alpha_+(y)=y^{-\sigma}e^{ty/2}\Gamma(y+\sigma+1)/(\Gamma(y+1)\lambda_+^{y+\sigma})$ and
$\alpha_-(y)=y^{\sigma}e^{-ty/2}\Gamma(y-\sigma+1)/(\Gamma(y+1)\lambda_-^{y-\sigma})$; $\alpha_-\ge e^{-\sigma t/2}$. Put
$\mathcal A_1=\sup_{\mathrm{Win}}|\alpha_+-\alpha_-|$ and $\mathcal A_2=\sup_{\mathrm{Win}}\alpha_-$; both $\mathcal A_1$ and $\mathcal A_2-1$
are $O(\Lambda_0^2/A)$.

\emph{Step 3 (Cauchy--Schwarz with (LM1)).} With $r_{\min}=\min_{\mathrm{Win}}r=r(\lfloor A\rfloor)$ one has $|\sum_{\mathrm{Win}}gC_m|\le(E_2/(\mathcal L(1,\ell)r_{\min}))^{1/2}\|\mathcal Q-\mathcal P\|_{\ell^2(\mathrm{Win})}$ and
$|\tilde T|\le(E_2/\mathcal L(1,\ell))^{1/2}\tau_T$, where $\tau_T^2=\sum_{y\le A}\chi T_A^2/r$.

\emph{Step 4 (three terms).} On $\mathrm{Win}$ the difference is $\mathcal Q-\mathcal P=(\alpha_+-\alpha_-)H_+(y+\sigma)+\alpha_-(H_+-H_-)(y+\sigma)+\alpha_-[H_-(y+\sigma)-H_-(y-\sigma)]$, where
$H_\pm=H_{\lambda_\pm}$. With $K_1=A+\sigma$ and the Dirichlet form $S_\lambda=\sum_k(k+1)(\Delta H_\lambda(k))^2$,
\[
\|\mathcal Q-\mathcal P\|_{\mathrm{Win}}\le\mathcal A_1\|H_{\lambda_+}\|
+\mathcal A_2\,t\sup_{\lambda\in[\lambda_-,\lambda_+]}\lambda^{-1}\bigl[((K_1+1)S_\lambda)^{1/2}+\|H_\lambda\|\bigr]
+\mathcal A_2\,2\sigma\bigl(S_{\lambda_-}/(x_L-\sigma+1)\bigr)^{1/2};
\]
the middle term comes from $\frac{d}{d\lambda}H_\lambda(k)=-\lambda^{-1}[(k+1)\Delta H_\lambda(k)+H_\lambda(k)]$ and Minkowski's
inequality over $\lambda$, the last from writing $H(y+\sigma)-H(y-\sigma)$ as a sum of $2\sigma$ differences, each index
$j\ge x_L-\sigma$ being used at most $2\sigma$ times.

\emph{Step 5 (one global bandwidth).} Lemma~\ref{lem:C-hankel} gives, with $\Delta_\lambda=1/\lambda_--1/\lambda_+$ and
$h_0=\|H_{\lambda_-}\|$, for $\lambda\in[\lambda_-,\lambda_+]$:
$\|H_\lambda\|\le e^{\Delta_\lambda B^2/2}h_0+Ne^{-B^2/(2\lambda_+)}$ and
$S_\lambda^{1/2}\le(B/\lambda_-^{1/2})\|H_\lambda\|+Ne^{-B^2/(2\lambda_+)}(B^2/\lambda_-+1)^{1/2}$.

\emph{Step 6 (back to $\mathcal P$).} By Step 2, $H_{\lambda_-}(k)=\mathcal P(k+\sigma)/\alpha_-(k+\sigma)$, so $h_0\le e^{\sigma t/2}\|\mathcal P\|_{\ell^2(\mathbb N_0)}$.
On $\mathrm{Win}$, (LM2) and the decrease of $r$ give $\chi(m-r)^2\le\mathcal L(1,M)\,r(x_L)\,\chi(m-r)g$, so
$\|\mathcal P\|\le(\mathcal L(1,M)r(x_L)E_2)^{1/2}+P_L+\tau'_T+\|F\|_{\rm out}$ with $P_L^2=\sum_{y<x_L}\chi(m-r)^2$ and $\tau'^2_T=\sum_{y\le A}\chi T_A^2$, and with $\|F\|^2_{\rm out}=\sum_{y>A}\chi F^2$.

\emph{Step 7 (closing).} Steps 4--6 give $\|\mathcal Q-\mathcal P\|_{\mathrm{Win}}\le\delta_1h_0+\rho_N$ with
\begin{align*}
\delta_1&=e^{\Delta_\lambda B^2/2}\Bigl[\mathcal A_1+\mathcal A_2\frac{t}{\lambda_-}\Bigl(\frac{(K_1+1)^{1/2}B}{\lambda_-^{1/2}}+1\Bigr)
+\mathcal A_2\frac{2\sigma B}{(\lambda_-(x_L-\sigma+1))^{1/2}}\Bigr],\\
\rho_N&=Ne^{-B^2/(2\lambda_+)}\Bigl[\mathcal A_1+\mathcal A_2\frac{t}{\lambda_-}\Bigl((K_1+1)^{1/2}\Bigl(\frac{B}{\lambda_-^{1/2}}+\Bigl(\frac{B^2}{\lambda_-}+1\Bigr)^{1/2}\Bigr)+1\Bigr)\\
&\qquad+\mathcal A_2\frac{2\sigma}{(x_L-\sigma+1)^{1/2}}\Bigl(\frac{B}{\lambda_-^{1/2}}+\Bigl(\frac{B^2}{\lambda_-}+1\Bigr)^{1/2}\Bigr)\Bigr].
\end{align*}
With Steps 1, 3 and 6, $E_2\le\delta'E_2+b\,E_2^{1/2}+\epsilon_{\rm out}$, where
$\delta'=\delta_1e^{\sigma t/2}(\mathcal L(1,M)r(x_L)/(\mathcal L(1,\ell)r_{\min}))^{1/2}$,
$b=[\delta_1e^{\sigma t/2}(P_L+\tau'_T+\|F\|_{\rm out})+\rho_N]/(\mathcal L(1,\ell)r_{\min})^{1/2}+\tau_T/\mathcal L(1,\ell)^{1/2}$ and
$\epsilon_{\rm out}=|e_L|+|e_{\rm out}|+|d_0|$. If $\delta'<1$ then
$E_2\le\bigl[(b+(b^2+4(1-\delta')\epsilon_{\rm out})^{1/2})/(2(1-\delta'))\bigr]^2$.

\emph{Asymptotics.} One has $Ne^{-B^2/(2\lambda_+)}\le2A^{-2}$; $\mathcal A_1$, $\mathcal A_2-1$, $\sigma t$, $t^2A$, $\beta_A^2A$, $\beta_L^2A$ and
$\Delta_\lambda B^2$ are $O(\Lambda_0^2/A)\to0$; $\delta_1=(1+o(1))\,tA^{1/2}B\,(1+\xi\kappa_L^{-1/2})$ with
$tA^{1/2}B=c_0^{3/2}/\phi_\xi\,(1+o(1))$; and $r(x_L)/r_{\min}\to\kappa_L^{-1/2}$ by the bounds on $r$ in Section~\ref{sec:C-range}.
Hence $\delta'\to c_0^{3/2}\Gamma_\xi<1$. With Lemma~\ref{lem:C-edge}, $b=O(A^{1/4})$ and $\epsilon_{\rm out}=O(A^{1/2})$, so
$E_2=O(A^{1/2})$.
\end{proof}

\begin{lemma}[Hankel form of the $\ell^2$ norms]\label{lem:C-hankel}
Let $J_0$ and $I_0$ be the Bessel function and the modified Bessel function of order zero. For a finite signed measure $\nu$ on
$[0,\infty)$ let $\hat\nu(s)=\int J_0(2sx^{1/2})\,\nu(dx)$. Then
\[
\sum_kp(k\,|\,\lambda x)p(k\,|\,\lambda x')=e^{-\lambda(x+x')}I_0\bigl(2\lambda(xx')^{1/2}\bigr)
=\frac2\lambda\int_0^\infty e^{-s^2/\lambda}J_0(2sx^{1/2})J_0(2sx'^{1/2})\,s\,ds,
\]
so that $\|H_\lambda\|^2=\frac2\lambda\int e^{-s^2/\lambda}|\hat\nu(s)|^2s\,ds$ and $S_\lambda=\frac2{\lambda^2}\int e^{-s^2/\lambda}|\hat\nu(s)|^2s^3\,ds$.
The bounds of Step 5 of Proposition~\ref{thm:C-decay} follow by splitting at $s=B$ and using $|\hat\nu|\le N$ above $B$.
\end{lemma}
\begin{proof}
The first equality is the series of $I_0$. The second is Weber's second exponential integral~\cite[13.31(1)]{Watson1922},
in its case of order zero, $\int_0^\infty e^{-p^2t^2}J_0(at)J_0(bt)\,t\,dt=(2p^2)^{-1}e^{-(a^2+b^2)/(4p^2)}I_0(ab/(2p^2))$ for
$|\arg p|<\pi/4$, with $a=2x^{1/2}$, $b=2x'^{1/2}$ and $p^2=1/\lambda$. The norm follows by integrating against
$\nu\otimes\nu$. For $S_\lambda$, compute $\frac{d}{d\lambda}(\lambda\|H_\lambda\|^2)$ on both sides: on the transform side, and as
$\sum_k(k+1)(\Delta H_\lambda(k))^2$ on the output side.
\end{proof}

\begin{lemma}[edge terms]\label{lem:C-edge}
With $\lambda_*=\max(\lambda_M,\lambda_\ell)$:
(i) $|d_0|\le(1+e)c_*+2e^{\Lambda_0-x_L}$;
(ii) $\tau_T^2\le c_*A^{-1/2}(1+(2\pi A)^{1/2})e^{\beta_A^2A/2}$ and $\tau'^2_T\le c_*\tau_T^2$;
(iii) $|e_{\rm out}|\le2e^{\beta_A^2A/2}(\lambda_*+2\lambda_*(A^{1/2}+1)+3)$;
(iv) $|e_L|\le\lambda_*\bigl[2+2e^{\beta_L^2A/2}+(1+M)c_*(1+2e^{1/4}(ex_L/(2\pi))^{1/2})\bigr]$;
(v) $P_L^2\le(1+M)^2c_*^2(1+e^{1/2}(e/(2\pi))(1+\log x_L))$;
(vi) $\|F\|^2_{\rm out}\le4(2\pi\lfloor A\rfloor)^{-1/2}$;
(vii) $c_*\ge(1-(e/4)^A)/(M(1+2e^{1/4}(eA/\pi)^{1/2}))$, hence $r_{\min}=\Omega(1/A)$.
\end{lemma}
\begin{proof}
(i) $\chi\le1$ on $[0,x_L)$ and $\delta_0(x)\le e^{-x}$ there; $\int e^{-x}d\rho_\infty=\mathsf P\rho_\infty(0)=r(0)=c_*$ and
$\sum_jw_je^{-x_j}=m(0)=e\,c_*$ (Proposition~\ref{prop:C-kkt}); on $[x_L,A]$ one has $\chi\delta_0\le e^{\Lambda_0-x_L}$, and the masses of
$\rho_A$ and $\pi$ are at most $1$.
(ii) $\chi T_A^2/r\le\chi T_A$ because $T_A\le r$; $T_A(y)\le c_*A^{-1/2}e^{-(A-y)^2/(2A)}$ and $\chi(y)\le e^{\beta_A(A-y)}$; sum over
$y\le A$ by the maximum plus the integral of the unimodal summand. Finally $\tau'^2_T=\sum\chi(T_A^2/r)r\le(\max r)\tau_T^2=c_*\tau_T^2$,
since $r$ is decreasing with $r(0)=c_*$.
(iii) For $y>A$ one has $|g(y)|\le\gamma_g(y):=\lambda_*+(y-y_A)(2\lambda_*+(y-y_A)/y_A)$ with $y_A=\lfloor A\rfloor$: the tangent of the convex
$G$ at $y_A$ has slope $G(y_A)-G(y_A-1)\ge\psi(y_A)-\lambda_\ell-\lambda_M$, and both $\varphi(y)-\varphi(y_A)\le(y-y_A)\psi(y)$ and $\psi(y)-\psi(y_A)\le(y-y_A)/y_A$ hold. We also have $\mathsf P(\chi|\mu|)(y)\le2e^{\beta_A^2A/2}p(y\,|\,A)$ for $y>A$, from
$p(y\,|\,x)\le p(y\,|\,A)e^{-(A-x)^2/(2A)}$ for $x\le A<y$ and $\chi(x)\le e^{\beta_A(A-x)}$. Then
$|e_{\rm out}|\le2e^{\beta_A^2A/2}E\gamma_g(Y')$ with $Y'\sim\mathrm{Poisson}(A)$, which gives (iii).
(iv) $\sum_{y<x_L}\mathsf P(\chi|\mu|)(y)=\int\chi(x)P_x(Y<x_L)\,d(\rho+\pi)$; for $x\ge x_L$ the Poisson Chernoff bound gives
$\chi(x)P_x(Y<x_L)\le e^{\beta_L(x-x_L)-(x-x_L)^2/(2x)}\le e^{\beta_L^2A/2}$; on $y\le A$ one has $|F|\le(1+M)r$, and $r$ is summed with the
bounds of Section~\ref{sec:C-range}.
(v) $P_L^2=\sum_{y<x_L}\chi(m-r)^2\le(1+M)^2\sum_{y<x_L}r^2$, with the bounds on $r$.
(vi) For $y>A$ one has $|F|\le R_A+m\le2p(y\,|\,A)$ because $p(y\,|\,\cdot)$ increases on $[0,y]$, and
$\max_yp(y\,|\,A)\le p(y_A\,|\,y_A)\le(2\pi y_A)^{-1/2}$.
(vii) $1=\sum m\le M\sum r$, and $\sum_yr(y)\le c_*(1+2e^{1/4}(eA/\pi)^{1/2})+(e/4)^A$ by the bounds on $r$, the last term bounding
the mass beyond $2A$.
\end{proof}
The constants in (i)--(vii) are finite powers and exponentials of $A^{1/3}$; they enter $b$ and $\epsilon_{\rm out}$ as $O(A^{1/4})$
and $O(A^{1/2})$ because $\beta_A^2A$ and $\beta_L^2A$ are $O(\Lambda_0^2/A)\to0$.

\begin{corollary}[one window]\label{cor:C-window}
Under Proposition~\ref{thm:C-decay}, for $A\ge A_0$ and every set $Y_W$ of integers inside $\mathrm{Win}$ one has $\sum_{y\in Y_W}(m-r)^2/r\le\mathcal L(1,M)\,C_0A^{1/2}\exp(-\Lambda_0\min_{Y_W}h(y/A))$, uniformly in $Y_W$.
\end{corollary}
\begin{proof}
(LM2) gives $(m-r)^2/r\le\mathcal L(1,M)(m-r)g$ on $\mathbb N_0$; each term of $E_2$ is nonnegative and
$\chi\ge\exp(\Lambda_0\min_{Y_W}h)$ on $Y_W$.
\end{proof}

\subsection{Counting mass points in a window}

\emph{Test functions.} For $\vartheta\in(0,1)$ and integers $1\le s_1<s_2$, the window
$\varpi(x)=P(s_1\le\mathrm{Poisson}(\vartheta x)\le s_2)$ has the exact deconvolution
$\Psi_\omega(y)=\sum_{s=s_1}^{s_2}\binom ys\vartheta^s(1-\vartheta+i\omega)^{y-s}$: $E_x\Psi_\omega(Y)=e^{i\omega x}\varpi(x)$, by binomial
thinning ($E_x[\binom Ys\vartheta^s(1-\vartheta+i\omega)^{Y-s}]=e^{i\omega x}e^{-\vartheta x}(\vartheta x)^s/s!$). We also have
$|\Psi_\omega(y)|\le\Psi_0(y)(1+\omega^2/(1-\vartheta)^2)^{(y-s_1)_+/2}$ with $\Psi_0(y)=P(s_1\le\mathrm{Bin}(y,\vartheta)\le s_2)\le1$.

\begin{lemma}\label{lem:C-count}
Let $l_1=s_1/\vartheta$, $l_2=s_2/\vartheta$, $\omega_0=2\pi/(l_2-l_1)$ and $W=[a_1,a_2]$ with $a_1\le l_1<l_2\le a_2\le A$. Let $k$ be the number of mass points of $\pi$ in $W$ and let $Y_W=[b_1,b_2]$ be a set of integers. Put $\epsilon_W=\max(\varpi(a_1),\varpi(a_2))$,
$\gamma=(l_2-l_1)\min_{[l_1,l_2]}(\varpi c_*\hat\rho)$, $\Omega_n=n\omega_0$ and $g_n(y)=(1+\Omega_n^2/(1-\vartheta)^2)^{(y-s_1)_+/2}$; put also $\tau_n=M\sum_{y\notin Y_W}g_n(y)\Psi_0(y)r(y)$, and let $\eta_W^2=\sum_{Y_W}(R_A-m)^2/r$ and $r_W=\sum_{Y_W}r$. Fix $r_b>1$, put
$n(k)=\lceil r_bk/(r_b-1)\rceil-1$, and let $n_*\ge1$ be an integer with
\[
\text{(a)}\ \ r_b\epsilon_W/\gamma\le1/4,\qquad\text{(b)}\ \ (r_b(2n_*+1))^{1/2}\tau_{n_*}/\gamma\le1/4 .
\]
If
$k_C:=\frac{r_b-1}{r_b}\,\frac{(1-\vartheta)(l_2-l_1)}{2\pi}\Bigl(\frac{2\log\bigl(\gamma/(2(r_b(2n_*+1)r_W)^{1/2}\eta_W)\bigr)}{b_2-s_1}\Bigr)^{1/2}>0$
and $n(\lceil k_C\rceil-1)\le n_*$, then $k\ge k_C$.
\end{lemma}
\begin{proof}
Suppose $k<k_C$ and put $n=n(k)$ ($n=0$ if $k=0$); then $n\le n(\lceil k_C\rceil-1)\le n_*$, so (b) holds at $n$ since $\tau_n$ increases
with $n$. Let $V$ be the space of trigonometric polynomials $\sum_{l=0}^nc_le^{il\omega_0x}$ vanishing at the $k$ mass points in $W$;
$\dim V\ge n+1-k$, so $(n+1)/\dim V\le r_b$. Take an orthonormal basis $(P_i)$ of $V$ in $L^2(\mu_\varpi)$, where $\mu_\varpi=\varpi\rho_A\,dx$, and
$Q_V=\sum|P_i|^2\ge0$; $Q_V$ vanishes at the mass points in $W$ and $\int Q_V\,d\mu_\varpi=\dim V$. The Gram matrix
$(\int e^{i(l'-l)\omega_0x}d\mu_\varpi)_{ll'}$ is at least $\gamma\,\mathrm{Id}$, with $\mathrm{Id}$ the identity matrix, because $\varpi c_*\hat\rho\ge\gamma/(l_2-l_1)$ on $[l_1,l_2]$, an
interval of length $2\pi/\omega_0$; so $Q_V\le(n+1)/\gamma$. Writing $Q_V=\sum_{|j|\le n}q_je^{ij\omega_0x}$ with coefficients $q_j$, the coefficient matrix,
conjugated by the square root of the Gram matrix, is a projector of rank $\dim V$; hence $\sum|q_j|^2\le(n+1)\dim V/\gamma^2$.

With $\Psi=\sum_jq_j\Psi_{j\omega_0}$, every finite measure $\nu'$ on $[0,A]$ satisfies $\int Q_V\varpi\,d\nu'=\sum_y\Psi\,(\mathsf P\nu')$; with
$\nu'=\rho_A-\pi$, $\dim V=\int Q_V\varpi\,d\pi+\sum_y\Psi(R_A-m)$, with absolute convergence because $|\Psi|$ grows at most geometrically
while $\mathsf P\nu'$ decays faster than any exponential. Now $\int Q_V\varpi\,d\pi\le(n+1)\epsilon_W/\gamma$, since $Q_V=0$ at the mass points
in $W$ and $\varpi\le\epsilon_W$ off $W$; $\sum_{Y_W}|\Psi_{j\omega_0}||R_A-m|\le g_n(b_2)r_W^{1/2}\eta_W$; off $Y_W$ the bound $|R_A-m|\le Mr$ holds, so
that part is at most $\tau_n$; and $\sum_j|q_j|v_j\le(\sum|q_j|^2)^{1/2}(2n+1)^{1/2}\max_jv_j$ for any $v_j\ge 0$. Dividing by $\dim V$ and using
$(n+1)/\dim V\le r_b$ gives $1\le r_b\epsilon_W/\gamma+(r_b(2n+1))^{1/2}[g_n(b_2)r_W^{1/2}\eta_W+\tau_n]/\gamma$, so by (a) and (b),
$g_n(b_2)\eta_W\ge\gamma/(2(r_b(2n+1)r_W)^{1/2})\ge\gamma/(2(r_b(2n_*+1)r_W)^{1/2})$. With $\log(1+z)\le z$ this gives $(n\omega_0)^2(b_2-s_1)/(2(1-\vartheta)^2)\ge\log\bigl(\gamma/(2(r_b(2n_*+1)r_W)^{1/2}\eta_W)\bigr)$, that is,
$n\ge\frac{(1-\vartheta)(l_2-l_1)}{2\pi}\bigl(2\log(\cdots)/(b_2-s_1)\bigr)^{1/2}$. As $n(k)<r_bk/(r_b-1)$, it follows that $k>\frac{r_b-1}{r_b}n\ge k_C$, a contradiction. For $k=0$, $n=0$ and $V=\mathbb C$ give
$1\le\epsilon_W/\gamma+[r_W^{1/2}\eta_W+\tau_0]/\gamma$, against (a), (b) and $k_C>0$.
\end{proof}

\subsection{Assembly of the proof}\label{app:C-proof}

Fix $\xi\in(0,1)$ and $c_0<c_{0}^{\max}(\xi)$, take $\delta\in(0,(1-\kappa_L)/2)$ and $r_b>2$, choose $\vartheta\in(0,1/2)$, and let $A\to\infty$. Put
$L_c=A^{3/4}$, the margin around the mass points $M_a=u(A/\vartheta)^{1/2}$ with $u^2=8\log A$, and the output margin
$M_o=2(2A\Lambda_0/\vartheta)^{1/2}$; both are $o(L_c)$, since $M_a=O((A\log A)^{1/2})$ and $M_o=O(A^{2/3})$. Place cores $[l_1,l_2]$ of
length $L_c$ in $[(\kappa_L+\delta)A,(1-\delta)A]$ with consecutive cores $2M_a+1$ apart ($s_1=\lceil\vartheta l_1\rceil$,
$s_2=\lfloor\vartheta l_2\rfloor$). Each core gets the window $W=[l_1-M_a,l_2+M_a]$; these windows are pairwise disjoint. Its output
window is $Y_W=[b_1,b_2]=[l_1-M_o,l_2+M_o]$; output windows may overlap, since each lemma is applied to one window. Let $h_W$ be the
minimum of $h$ over $Y_W/A$ and put
$\bar\eta_W^2=2\mathcal L(1,M)C_0A^{1/2}e^{-\Lambda_0h_W}+2\sum_{y\in Y_W}T_A^2/r$. By Corollary~\ref{cor:C-window}, $\bar\eta_W\ge\eta_W$, and
Lemma~\ref{lem:C-count} holds with $\bar\eta_W$ in place of $\eta_W$, because its proof gives
$g_n(b_2)\eta_W\ge\gamma/(2(r_b(2n_*+1)r_W)^{1/2})$. Uniformly over the windows:
\begin{enumerate}
\item The inequalities $T_A\le c_*A^{-1/2}e^{-(A-y)^2/(2A)}$ and $y\le(1-\delta)A+M_o$ bound the second term of $\bar\eta_W^2$ by $e^{-\delta^2A/3}$,
so $\log(1/\bar\eta_W)=(\Lambda_0/2)h_W-O(\log A)$, with $h_W\ge h_\delta:=\min(h(\kappa_L+\delta/2),h(1-\delta/2))>0$ for large $A$. In
particular $cA^{1/3}\le\log(1/\bar\eta_W)\le\Lambda_0$.
\item The estimates $c_*=\Theta(A^{-1/2})$ and $\hat\rho(x)\ge c_1/(x\log x)^{1/2}$ (Lemma~\ref{lem:C-sizes}) and $\varpi\ge1/2-o(1)$ on the core give
$\gamma\ge c'L_cA^{-1/2}(A\log A)^{-1/2}=c'A^{-1/4}(\log A)^{-1/2}$. The Poisson Chernoff bounds at the two ends of $W$ give
$\epsilon_W\le\exp(-u^2l/(2(l+M_a))(1-o(1)))\le A^{-3}$ for large $A$. So (a) holds.
\item The count of Lemma~\ref{lem:C-count} on a window is $f_W=\Theta(L_c(\log(1/\bar\eta_W)/l_2)^{1/2})=\Theta(A^{5/12})$; take
$n_*=n(\lceil f_W\rceil)=\Theta(A^{5/12})$, so $\Omega:=n_*\omega_0=\Theta(A^{-1/3})$. For $y>b_2$ one has $\Psi_0(y)\le P(\mathrm{Bin}(y,\vartheta)\le s_2)\le\exp(-(\vartheta y-s_2)^2/(2\vartheta y))$ and
$g_{n_*}(y)\le\exp(\Omega^2(y-s_1)/(2(1-\vartheta)^2))$. At $y=b_2$ the exponent satisfies $(\vartheta y-s_2)^2/(2\vartheta y)\ge\vartheta M_o^2/(2b_2)\ge4\Lambda_0$,
while $\Omega^2b_2/(2(1-\vartheta)^2)\le(1+o(1))\log(1/\bar\eta_W)/(1-\vartheta)\le2.1\Lambda_0$ (as $\vartheta<1/2$). Beyond $b_2$ the
derivative of the tail exponent is at most $-(\vartheta/2)(1-(s_2/(\vartheta b_2))^2)=-\Theta(A^{-1/3})$, against the growth rate
$\Omega^2/(2(1-\vartheta)^2)=\Theta(A^{-2/3})$. For $y<b_1$ one has $P(\mathrm{Bin}(y,\vartheta)\ge s_1)\le\exp(-(s_1-\vartheta y)^2/(2(\vartheta y+(s_1-\vartheta y)/3)))\le e^{-3\Lambda_0}$ at $y=b_1$ and
smaller below, while $g_{n_*}(y)\le g_{n_*}(b_1)\le e^{2.1\Lambda_0}$. With $r\le c_*$ one gets $\tau_{n_*}\le e^{-\Lambda_0/2}$ for large $A$, so (b)
holds with $r_b>2$ fixed, and $\log\bigl(\gamma/(2(r_b(2n_*+1)r_W)^{1/2}\bar\eta_W)\bigr)=\log(1/\bar\eta_W)-O(\log A)$.
\item One has $b_2-s_1=l_2(1-\vartheta)(1+O(A^{-1/4}))$ and $l_2=x(1+O(A^{-1/4}))$ for $x$ in the core.
\end{enumerate}
So Lemma~\ref{lem:C-count} gives, on each window,
$k_W\ge(1-o(1))(1-1/r_b)(1-\vartheta)^{1/2}\frac{l_2-l_1}{2\pi}(\Lambda_0h_W/l_2)^{1/2}$, and $h_W=h(x/A)(1-o(1))$ for $x$ in the
core, since $h$ is continuous and positive on $[\kappa_L+\delta,1-\delta]$ and the window has relative length $o(1)$. The windows are
disjoint, so $K(A)\ge\sum_Wk_W$. The cores cover $[(\kappa_L+\delta)A,(1-\delta)A]$ except the spacings, of total length
$O(A^{1/4}M_a)=O(A^{3/4}(\log A)^{1/2})=o(A)$; dropping them from the integral below loses $O(A^{3/4}(\log A)^{1/2}A^{-1/3})=o(A^{2/3})$.
Summing as a Riemann sum,
\begin{align*}
K(A)&\ge(1-o(1))(1-1/r_b)(1-\vartheta)^{1/2}\frac{\Lambda_0^{1/2}}{2\pi}\int_{(\kappa_L+\delta)A}^{(1-\delta)A}\Bigl(\frac{h(x/A)}{x}\Bigr)^{1/2}dx\\
&=(1-o(1))(1-1/r_b)(1-\vartheta)^{1/2}\frac{c_0^{1/2}}{2\pi}A^{2/3}\int_{\kappa_L+\delta}^{1-\delta}\Bigl(\frac{h(f)}{f}\Bigr)^{1/2}df .
\end{align*}
Divide by $A^{2/3}$, take the $\liminf$, then let $\delta$ tend to $0$, $r_b$ to $\infty$, $\vartheta$ to $0$ and $c_0\to c_{0}^{\max}(\xi)$, and maximize
over $\xi$. The value of $\lambda_M$, the bounds of Proposition~\ref{prop:C-lower}(a) and the constant are enclosed in ball
arithmetic. At $\xi=0.2411$ (an exact decimal), $\kappa_L$ is isolated by bisection with signs proved by interval arithmetic, and
$c_{0}^{\max}\in[0.06490794803,0.06490794804]$. The integral $\mathcal I$ is bounded below by a rigorous quadrature over $[\kappa_L+10^{-8},1-10^{-8}]$, where the integrand is analytic, and above by
adding the two end pieces with the bound $\kappa_L^{-1/2}$ on the integrand; this gives $\mathcal I\in[1.3824928,1.3824930]$ and
$c_0^{\max}(\xi)^{1/2}\mathcal I(\xi)/(2\pi)\in[0.0560572,0.0560573]$. Since $c_P$ is a maximum over $\xi$, it follows that $c_P>0.0560$. \qed

\subsection{A spontaneous count}\label{app:C-lam}

Fix $\lambda>0$, let $\pi$ be the capacity-achieving input of $Y\sim\mathrm{Poisson}(X)$ with $\lambda\le X\le U:=\lambda+A$, and let
$C$ be its capacity. Fact (F1) holds for this channel, with $[\lambda,U]$ in place of $[0,A]$~\cite{Shamai1990}~\cite[Lemma~1]{BDS2024};
(F3) and (F4) hold for every input law~\cite[Lemma~3]{BDS2024}; (F2) is not used below. Define $r$, $c_*$ and $g$ as before, with this
$C$. Lemma~\ref{lem:C-digamma} and the identity of Proposition~\ref{prop:C-kkt} do not depend on the input range, so the optimality
conditions read $E_xg(Y)\ge\delta_0(x)$ on $[\lambda,U]$, with equality at every mass point. The value $g(0)=1$ is not available. It was
used in Proposition~\ref{prop:C-upper} when $g$ is maximal at $0$, in Proposition~\ref{prop:C-lower} at $y_1=0$, and in
Lemma~\ref{lem:C-edge}(i) through $m(0)=e\,c_*$. The lemma below replaces the first use.

\begin{lemma}[the value at zero]\label{lem:C-zero}
For every $\lambda\ge0$, with $m_0=E[X\,|\,Y=0]$ and $m_1=E[X\,|\,Y=1]$, one has $g(0)\le1-m_0+m_0\log(m_1/m_0)$. If $g$ is maximal
at $0$, then $g(0)\le1+2e^{-1-\gamma_E}<1.42<\lambda_M$.
\end{lemma}
\begin{proof}
By (F3), $G(1)-G(0)=\log m_0$, so convexity of $G$ gives $G(y)\ge G(0)+y\log m_0$ on $\mathbb N_0$; with $g=G-\varphi+C+1$ and
$\varphi(0)=0$ this gives $g(y)\ge g(0)+y\log m_0-\varphi(y)$. Take $E_{x_j}$ at a mass point $x_j$; as in Lemma~\ref{lem:C-pointwise},
$\delta_0(x_j)$ cancels and $g(0)\le1-m_0+B(x_j,m_0)$. Average over the posterior of $X$ given $Y=0$, which has mean $m_0$. The posterior
given $Y=1$ is that posterior tilted by $x$, so $E[B(X,m_0)\,|\,Y=0]=m_0E[\log(X/m_0)\,|\,Y=1]\le m_0\log(m_1/m_0)$ by Jensen's
inequality. If $g$ is maximal at $0$, then $g(2)\le g(0)$; since $G(2)-G(0)=\log m_0+\log m_1$ and $\varphi(2)=\psi(1)+\psi(2)=1-2\gamma_E$,
this gives $\log m_1\le1-2\gamma_E-\log m_0$. Hence $g(0)\le1-2\gamma_Em_0-2m_0\log m_0$, and the right side is largest at
$m_0=e^{-1-\gamma_E}$, where it equals $1+2e^{-1-\gamma_E}=1.4131\ldots$.
\end{proof}

With this lemma the proof of Proposition~\ref{prop:C-upper} holds for $\lambda>0$, since its case $y_0\ge1$ uses only the mass points,
the posteriors and the values of $g$ at $y_0\pm1$; so $m\le Mr$ on $\mathbb N_0$. Proposition~\ref{prop:C-lower} holds for every integer
$y_1\in[\lambda,U]$: its proof uses the optimality inequality at the feasible point $x=y_1$, the case $y_1=0$ does not arise, and the
bounds of its part (a) do not depend on the channel. Corollary~\ref{cor:C-LM} follows with $[\lambda,U]$ in place of $[0,A]$ in (LM1).
For the integers $y<\lambda$, put $y_1=\lceil\lambda\rceil$; convexity of $G$ with the slope $\log E[X\,|\,Y=y_1-1]\le\log U$ gives
$g(y)\ge-\lambda_\ell-y_1(\gamma_E+\log U)$, so $|g(y)|=O(\log A)$ there.

Take $\rho$ to be the restriction of $\rho_\infty$ to $[\lambda,U]$; Lemma~\ref{lem:C-ref}(iv) holds for it, since its proof uses only
that $\rho/\rho([\lambda,U])$ is a feasible input. Then $\mathsf P\rho=r-T_U-L_\lambda$ on $\mathbb N_0$, where $T_U$ is $T_A$ with $U$ in place
of $A$ and $L_\lambda=\mathsf P(\rho_\infty\mathbf 1_{[0,\lambda)})$. As $p(y\,|\,x)e^x\le\lambda^y/y!$ for
$x\le\lambda$ and $\int e^{-x}\rho_\infty(dx)=r(0)=c_*$, one has $L_\lambda(y)\le c_*\lambda^y/y!$. The proof above now runs with $U$ in
place of $A$ in the weight $\chi$, the window and the cores, with three changes:
\begin{enumerate}
\item In Step~1 of Proposition~\ref{thm:C-decay}, $F=r-T_U-L_\lambda-m$, and the identity gains the term $-\tilde L$ with
$\tilde L=\sum_{y\le U}\chi gL_\lambda$. Below $x_L$ one has $\chi\le1$ and $|g|=O(\log A)$; on and above $x_L$ the weight satisfies $\chi\le e^{\Lambda_0}$, while
$\lambda^y/y!\le(e\lambda/y)^y$ decays faster than any exponential. So $|\tilde L|=O(c_*\log A)$, which enters $\epsilon_{\rm out}$.
The part of $\tilde T$ below $\lambda$ is bounded in the same way, since $\chi(y)\le(\lambda/U)^{2\sigma}e^{tU}$ there; the rest of
$\tilde T$ is bounded as in Step~3.
\item In Step~6, $\|\mathcal P\|$ gains the term $(\sum_y\chi L_\lambda^2)^{1/2}=O(c_*)$, which enters $b$; in Lemma~\ref{lem:C-count} and
the assembly, $T_U+L_\lambda$ replaces $T_A$ in $\bar\eta_W$, and $L_\lambda$ is negligible on the output windows.
\item In Lemma~\ref{lem:C-edge}(i), $m(0)=e^{g(0)}c_*\le Mc_*$, so $|d_0|\le(1+M)c_*+2e^{\Lambda_0-x_L}$. In (iv), $|F|\le Mr$ on
$y\le U$ because $0\le T_U+L_\lambda\le r$; in (iii), $y_A=\lfloor U\rfloor$ lies in $[\lambda,U]$ for $A\ge1$. Lemma~\ref{lem:C-sizes}(ii)
sums over the integers of $[\lambda,U]$ and still gives $c_*=\Theta(A^{-1/2})$.
\end{enumerate}
None of these changes alters $\delta'$, so its limit and the conditions $b=O(A^{1/4})$ and $\epsilon_{\rm out}=O(A^{1/2})$ stand. The
assembly gives $K(A,\lambda)\ge(1-o(1))(c_P-\varepsilon)\,U^{2/3}$ for every $\varepsilon>0$, and $U^{2/3}/A^{2/3}\to1$, which proves the
theorem for $\lambda>0$. \qed
